% \section{Conclusion}
\section{结论}

% In this paper, we propose \textbf{\model{}}, a framework that enables users to effectively interact with human "dreams" and influence final presentations.
在本文中，我们提出了 \textbf{\model{}}，这是一个能够让用户有效地与人类“梦境”进行交互并影响其最终呈现结果的框架。
% Our framework offers several appealing properties:
我们的框架具有以下几个吸引人的特性：
% 1) We \emph{pioneer a novel interaction paradigm with the human brain} by introducing a prototype that directly instructs fMRI signals via natural language descriptions, opening up exciting possibilities for future applications.
1) 我们\emph{开创了一种与人类大脑交互的新型范式}，通过引入一个原型系统，利用自然语言描述直接对 fMRI 信号发出指令，为未来的应用开辟了激动人心的可能性。
% 2) Our model's capability can be further enhanced using techniques that consider temporal and spatial perspectives, thereby improving the fidelity and precision of instructions.
2) 通过引入考虑时空视角的各项技术，我们可以进一步增强模型的能力，从而提高指令的保真度与精确度。
% 3) With minimal effort, our system can seamlessly extend to support multi-modal instructions.
3) 我们的系统只需极少量的改动，即可无缝扩展以支持多模态指令。

% \paragraph{Limitations and Future Work.} 1) In this work, we consider the obtained biological brain signals from visual stimuli as representations of human ``thoughts''. However, in real-world applications, the situation is more complex. ``Dreams,'' for instance, might originate from internal brain activity rather than external stimuli. 2) Our model struggles with instructions involving the addition of small objects. 3) For the LLM-enhanced region-aware instructions, an extra forward process is required to first extract the instruction-relevant regions from the intermediate reconstructions. 4) Future work will explore other types of ``dreams,'' such as audio and text, and more complex interaction strategies, such as multi-turn conversations.
\paragraph{局限性与未来工作。} 1) 在本研究中，我们将从视觉刺激中获取的生物脑信号视为人类“思想”的表征。然而在实际应用中，情况往往更加复杂。例如，“梦境”可能源自内部的大脑活动，而非外部刺激。2) 我们的模型在处理涉及添加细小对象的指令时存在困难。3) 对于基于大语言模型（LLM）增强的区域感知指令，需要额外的正向计算过程，才能从中间重建结果中提取出与指令相关的区域。4) 未来的工作将探索其他类型的“梦境”（如音频和文本），以及更复杂的交互策略（如多轮对话）。

% \paragraph{Ethical Issues.} As the title suggests, this model can be potentially exploited for malicious purposes such as mis-interpreting thoughts within human brain. We are committed to limiting the usage of our model strictly to research purposes.
\paragraph{伦理问题。} 正如标题所暗示的那样，该模型可能会被恶意利用，例如误解人类大脑中的真实想法。我们承诺严格限制我们模型的使用范围，确保其仅被用于研究目的。